\section{Network and Congestion Model}
\label{sec:network}

\subsection{Study area}

We study a $\sim$4$\times$4\,km area of central Bengaluru, India, covering
MG Road, Cubbon Park, Shivajinagar and Richmond Town. The area was chosen
for its irregular, organically grown topology, which is structurally unlike
the grid layouts and synthetic graphs common in this literature. The network
is obtained from OpenStreetMap via OSMnx~\cite{boeing2017osmnx} as a directed
multigraph, reduced to its largest strongly connected component so that every
origin--destination pair is reachable, giving \textbf{2{,}856 nodes and
7{,}015 directed edges}, with \textbf{mean out-degree 2.46}.

That last figure is the single most important structural property for what
follows. A node with out-degree two offers a walk-based extraction procedure
almost no choice, and one-way streets mean that a direction favoured by an
undirected flow computation may be illegal to traverse. Both conditions are
absent from the complete and near-complete undirected graphs used in prior
Physarum evaluations.

\subsection{Congestion as an interface, with two sources}

Algorithms never observe the congestion source directly. A
\texttt{Congestion\-Provider} returns a \texttt{TrafficSnapshot} mapping each
edge to a congested travel time, which keeps the choice of source orthogonal
to the routing code. We implement two providers.

\textbf{SUMO microsimulation (primary).} The same OpenStreetMap extract is
converted with \texttt{netconvert} and loaded with synthetic demand generated
by \texttt{randomTrips.py} and \texttt{duarouter} at two levels, light and
heavy. Per-edge travel times are collected from \texttt{edgeData} output at
300\,s intervals. We deliberately omit \texttt{--geometry.remove} during
conversion, because it merges edges and destroys the way-ID traceability the
join below depends on.

\textbf{BPR volume-delay (secondary).} An analytic volume-delay
function~\cite{bpr1964} with synthetic origin--destination demand and
incremental all-or-nothing assignment. It is retained because it is
independent of the simulator and because, as Section~\ref{sec:results}
shows, the two sources disagree in a way that is itself a result.

\subsection{Joining SUMO edges to graph edges}

SUMO identifiers derive from OSM way IDs, with sign encoding direction and a
\texttt{\#n} suffix marking split segments; stripping these recovers the way
ID and yields a mapping from each graph edge to its matched simulator
segments. We achieve \textbf{97.53\% ID-based coverage}, so no geometric
fallback matching was required; unmatched edges retain free-flow times and
are counted.

One detail here is worth stating because getting it wrong corrupts every
downstream number silently. \texttt{netconvert} splits ways while OSMnx
merges them, so the ratio of matched-segment length to graph-edge length has
median 1.44 and 90th percentile 8.9 on this network. Averaging the matched
segments' \emph{travel times} therefore answers ``how long to cross one
matched segment'', not ``how long to cross this graph edge'', and understates
each edge in proportion to how finely it was split, which affects busy arterials most. Under that rule 1{,}794 of 7{,}015 edges (26\%) implied
speeds above the fastest limit anywhere in the network, the worst by a factor
of 266, which silently violates A*'s admissibility. Aggregating in
\emph{pace} instead,
\begin{equation}
t_e \;=\; \ell_e \cdot \frac{\sum_i t_i}{\sum_i \ell_i},
\end{equation}
where $\ell_e$ is the graph edge length and $(t_i, \ell_i)$ are the matched
segments' travel times and lengths, reduces the count of over-maximum edges
to 1 and restores exact A*--Dijkstra agreement.

\subsection{A single shared free-flow baseline}

Stretch is only meaningful against a denominator on the same weighted graph.
We therefore re-annotate free-flow speeds from the generated SUMO network, so
that the congested numerator and the free-flow denominator derive from one
road model, and stamp the choice of baseline into every snapshot's metadata.
A consequence must be stated plainly: mean free-flow speed rises from
22.74 to 51.49\,km/h, so \textbf{absolute travel times on this graph are
optimistic for Bengaluru}. Ratios, which is all we report, are comparable; absolute seconds are not realistic.

\subsection{Calibration and what the two sources do and do not agree on}

At heavy demand, teleports (SUMO's escape hatch for gridlocked vehicles,
which falsify edge travel times) are held to 5.13\%. Per-interval travel
times are clipped at 20$\times$ free flow; this clip is necessary even in a
zero-teleport, fully drained run, and the clipped samples' median is
86--97$\times$ free flow, a cleanly separated tail rather than a threshold
judgement.

Measured on one denominator, congestion ratios along OD paths are
1.18$\times$ (light) and 2.53$\times$ (heavy) for BPR, against 1.81$\times$
and 2.11$\times$ for SUMO. \textbf{The two sources bracket rather than agree,
and must not be merged into a single notion of ``congestion level''.} SUMO's
floor is higher because signal delay enters the numerator but not the
free-flow denominator; BPR's range is wider because synthetic demand loads
shortest paths, which is exactly where an OD-path measurement walks. The
heavy-versus-light contrast is modest per edge (1.55$\times$ to
1.85$\times$): at light demand 24\% of graph edges carry no vehicle at all,
and 12\% at heavy.
